\workshoptitle{Bring your own laptop}{Windows 10 / 11}

\begin{context}
Download \textbf{rps-event-kit.zip} from the release page below. Right-click the
ZIP and choose Extract All. Find the folder containing \textbf{start.sh},
\textbf{wheelhouse} and \textbf{requirements.txt}; it may be one folder deeper
than the first folder Windows opens.
\end{context}

{\small\url{https://github.com/Eastbridge-Academy/rps-client-kit/releases/latest}\par}

\lessonstep{Step 1}{Install uv and Python}
Open PowerShell from the Start menu. If \code{uv --version} already works,
skip the installer command. Otherwise, run:

\begin{lstlisting}[style=plaincode]
winget install --id astral-sh.uv -e
\end{lstlisting}

Close PowerShell and open it again, then run:

\begin{lstlisting}[style=plaincode]
uv python install 3.12
\end{lstlisting}

If winget is unavailable, ask a facilitator or use the Windows installer at
\url{https://docs.astral.sh/uv/getting-started/installation/}.
These downloads need internet access; installing the kit and practicing use
local files.

\lessonstep{Step 2}{Create your bot project}
In File Explorer, open the extracted folder containing \textbf{requirements.txt}.
Click the address bar, type \code{powershell}, and press Enter. Run the commands
below; each final backtick continues the command onto the next line.

\begin{lstlisting}[style=plaincode]
$RpsProject = "$HOME\rps-bot"
uv venv --offline --python 3.12 "$RpsProject\.venv"
uv pip install --offline --no-index `
  --python "$RpsProject\.venv\Scripts\python.exe" `
  --find-links .\wheelhouse -r .\requirements.txt
cd $RpsProject
.\.venv\Scripts\rps-cli.exe init
.\.venv\Scripts\rps-cli.exe test
.\.venv\Scripts\rps-cli.exe validate
\end{lstlisting}

Open \textbf{bot.py} in the new \textbf{rps-bot} folder. In your editor, select
\textbf{.venv/Scripts/python.exe} as the Python interpreter.

\lessonstep{Step 3}{Use the handouts}
Where a handout says \code{rps-cli}, use
\code{.\textbackslash{}.venv\textbackslash{}Scripts\textbackslash{}rps-cli.exe}.
Each new PowerShell window needs \code{cd "\$HOME\textbackslash{}rps-bot"}
first. You do not need to activate a PowerShell script or change its execution policy.

\begin{checkpoint}
\textbf{Join the event:} follow \emph{Submit your bot} on page 1, using the
Windows executable path. Reuse the same team name for every upload.
\end{checkpoint}
